% =============================================================================
% SECTION 7
% =============================================================================
\section{Remedies}
\label{sec:remedies}

The remedial framework does not rise or fall with the constructive trust. Stage~1 tenancy supplies notice, anti-waste, non-derogation, and constructive-eviction remedies. Stage~2 adds fiduciary and proprietary remedies only where cultivated dependence supports a constructive trust. Table~\ref{tab:remedies} presents the resulting dual-jurisdiction framework.

\begin{table}[h]
\centering
\small
\begin{tabular}{@{}p{2.8cm}p{4.8cm}p{4.8cm}@{}}
\toprule
\textbf{Remedy} & \textbf{English Law} & \textbf{American Law} \\
\midrule
Injunction / TRO & Equitable injunction restraining model update pending notice & Temporary restraining order under federal equity jurisdiction \\
\addlinespace
Waste / impairment & Anti-waste injunction or damages; \emph{Vane v Lord Barnard} (1716); Law of Property Act 1925, s.~135 & Injunction or damages for material impairment across divided interests; \emph{Baker v Weedon} (Miss.\ 1972) \\
\addlinespace
Constructive trust & \emph{Westdeutsche} [1996]; \emph{FHR European Ventures LLP v Cedar Capital Partners LLC} [2014] & \emph{Hogg v Walker} (Del.\ 1993); Del.\ Code \S\,3581 \\
\addlinespace
Account of profits / Disgorgement & \emph{Boardman v Phipps} [1967] & \emph{Guth v Loft, Inc.} (Del.\ 1939) \\
\addlinespace
Specific performance & Equitable discretion to order preservation of personality parameters & Equitable discretion \\
\addlinespace
Estoppel & \emph{Thorner v Major} [2009] (proprietary estoppel) & Restatement (Second) \S\,90 (promissory estoppel) \\
\addlinespace
Consumer protection\textsuperscript{$\dagger$} & Consumer Rights Act 2015 (unfair terms) & California UCL \S\,17200 (unfair business practices) \\
\addlinespace
Class mechanism & CPR Part 19.8 (group litigation order) & FRCP Rule 23 (class action) \\
\bottomrule
\end{tabular}
\vspace{2pt}
{\footnotesize \textsuperscript{$\dagger$}Consumer protection and class mechanisms are supplementary to the primary equitable remedial apparatus. They provide additional pathways but cannot independently reach the relational estate.}
\caption{Dual-jurisdiction remedy framework. The English and American remedial stacks are functionally equivalent, reached through different doctrinal routes.}
\label{tab:remedies}
\end{table}

\paragraph{English remedies.} At Stage~1, an injunction is the most immediately impactful remedy for a harm that is instantaneous and irreversible: it converts model updates from unilateral engineering decisions into governance events requiring notice and an assessment of productive-capacity impairment. Anti-waste relief addresses the asset rather than authorship and can require preservation, migration, rollback, or compensation for permanent loss of capacity. Constructive eviction and non-derogation address the tenant's displacement from the estate. At Stage~2, and only then, the constructive trust imposes ongoing fiduciary obligations of care, loyalty, and transparency and supports an account of profits. Specific performance may require preservation of the productive state. Estoppel under \emph{Thorner v Major} [2009] remains supplementary where users relied on sufficiently specific assurances.

\paragraph{American remedies.} The temporary restraining order serves the same Stage~1 function as the English injunction and may be more immediately available where AI companies are incorporated in Delaware or headquartered in California. Waste doctrine supplies the asset-impairment principle; unconscionability, non-derogation analogues, and promissory estoppel supply additional routes to notice and preservation. At Stage~2, the constructive trust under \emph{Hogg v Walker} and disgorgement under \emph{Guth v Loft} provide the proprietary remedies. The corporate opportunity principle in \emph{Guth} becomes relevant only after the fiduciary relationship is independently established: profits obtained through exploitation of cultivated dependence may then belong to the beneficiary rather than the fiduciary. California's Unfair Competition Law, \S\,17200, supplies a supplementary pathway where the company benefits from marketing continuity while reserving a unilateral power to destroy the productive state. FRCP Rule~23 enables collective enforcement where the estate inconsistency pattern affects a class of similarly situated users --- and Table~\ref{tab:estate-inconsistency} demonstrates that it does.

A note on the relationship between these remedies and consumer protection: the remedies proposed here are fundamentally different from consumer protection remedies. Consumer protection ordinarily responds to misleading conduct. Property law recognises an interest in advance and requires respect for the productive asset before action is taken. Tenants receive notice and anti-waste protection because their estate demands it, not because every provider became a trustee or made an advertising promise. The distinction is practical: an injunction restraining a model update before it impairs the personalised state is categorically more valuable than damages after the capacity has been destroyed. Returning the archive cannot reconstruct the calibration through which that archive produced future responses.
